% !TeX root = ../thesis.tex
% Kapitel 2
\chapter{Grundlagen}\label{cha:grundlagen}

\section{Security Risk Management und Kontrollen}\label{sec:grundlagen-risk}
\subsection{Risiko und Risikomanagementprozess}
Einen allgemeinen Rahmen für den Umgang mit Risiken in Organisationen gibt die internationale Norm ISO~31000 vor~\cite{iso31000}.
Sie definiert ein \emph{Risiko} als \enquote{Auswirkung von Unsicherheit auf Ziele}~\cite[S.~7]{iso31000}.
Beschrieben wird ein Risiko üblicherweise über seine Ursachen, die möglichen Ereignisse, deren Auswirkungen und die Eintrittswahrscheinlichkeit~\cite[S.~7]{iso31000}.
Darauf aufbauend versteht die Norm unter \emph{Risikomanagement} die koordinierten Aktivitäten, mit denen eine Organisation in Bezug auf Risiken gelenkt und gesteuert wird~\cite[S.~7]{iso31000}.

Wie Risikomanagement konkret abläuft, beschreibt die Norm als Prozess~\cite[S.~16--23]{iso31000}.
Nach dem Festlegen von Anwendungsbereich, Kontext und Risikokriterien folgt die \emph{Risikobeurteilung}, die sich in Risikoidentifikation, Risikoanalyse und Risikobewertung gliedert.
Anschließend werden in der \emph{Risikobehandlung} geeignete Maßnahmen ausgewählt und umgesetzt.
Den gesamten Ablauf begleiten Kommunikation und Konsultation, Überwachung und Überprüfung sowie das Aufzeichnen und Berichten der Ergebnisse.
Auch wenn die Schritte meist nacheinander dargestellt werden, verläuft der Prozess in der Praxis iterativ~\cite[S.~17]{iso31000}.

ISO~31000 ist bewusst allgemein gehalten und auf jede Art von Risiko anwendbar.
Für den Bereich der Informationssicherheit konkretisiert die Norm ISO/IEC~27005 diesen Prozess und stimmt ihn mit den Anforderungen an ein Informationssicherheitsmanagementsystem nach ISO/IEC~27001 ab\cite[S.~7, 15]{iso27005}.
Das Management von Informationssicherheitsrisiken nach diesem Vorgehen wird im Folgenden als \emph{Security Risk Management} bezeichnet.

\subsection{Risk Register}\label{subsec:risk_register}
Sowohl ISO~31000 als auch ISO/IEC~27001 verlangen, dass die Ergebnisse der Risikobeurteilung und -behandlung dokumentiert werden~\cites[S.~23]{iso31000}[S.~10--11]{iso27001}.
Eine bestimmte Form der Dokumentation schreiben die Normen jedoch nicht vor.
In der Praxis werden diese Informationen häufig in einem \emph{Risk Register} festgehalten, einem zentralen Verzeichnis, in dem Risikoinformationen über alle Prozessschritte hinweg und über die Zeit gesammelt werden~\cite[S.~1]{nistir8286}.
Jedes Risiko wird darin als Eintrag mit festen Attributen erfasst.
Typisch sind eine Beschreibung, eine Kategorie, die Wahrscheinlichkeit und Auswirkung, die geplante Risikobehandlung, der Risikoeigentümer und der Status~\cite[S.~12--13]{nistir8286}.
Da das Register konsistent und iterativ genutzt wird und nach jeder Überprüfung ein neuer Zyklus beginnt~\cites[S.~13]{nistir8286}[S.~48]{iso27005}, enthält es neben aktuellen auch bereits bewertete Risiken mit ihren Maßnahmen.

\subsection{Kontrollen und Kontrollkataloge}\label{subsec:kontrollen}
Die Maßnahmen, mit denen ein Risiko behandelt wird, werden als \emph{Kontrollen} (engl.~\emph{controls}) bezeichnet.
ISO~31000 definiert eine Kontrolle als Maßnahme, die ein Risiko beibehält oder verändert.
Dazu zählen Prozesse, Grundsätze, Instrumente und Verfahren ebenso wie andere Bedingungen oder Aktionen mit dieser Wirkung~\cite[S.~9]{iso31000}.\footnote{Die deutschen Normfassungen übersetzen \emph{control} mit \enquote{Steuerung} bzw.~\enquote{Informationssicherheitsmaßnahme}~\cites[S.~9]{iso31000}[S.~18]{iso27001}. In dieser Arbeit wird der in der Praxis übliche Begriff \enquote{Kontrolle} verwendet.}
Für jedes Risiko, das behandelt werden muss, sollten eine oder mehrere Kontrollen festgelegt werden~\cite[S.~35]{iso27005}.

Als Orientierung dienen \emph{Kontrollkataloge}.
ISO/IEC~27001 enthält in Anhang~A 93 Kontrollen, die in vier Themen gegliedert sind: organisatorische, personenbezogene, physische und technologische Kontrollen~\cite[S.~18--26]{iso27001}.
Organisationen müssen die von ihnen festgelegten Kontrollen mit diesem Katalog abgleichen, um sicherzustellen, dass keine notwendige Kontrolle fehlt~\cite[S.~11]{iso27001}.
Der Katalog ist jedoch nicht abschließend und Organisationen können eigene Kontrollen definieren, deren Wortlaut möglichst genau beschreiben sollte, was die Kontrolle bewirkt~\cites[S.~11]{iso27001}[S.~35]{iso27005}.

Einen anderen Ansatz verfolgt das NIST Cybersecurity Framework~(CSF)~2.0.
Statt einzelne Kontrollen vorzugeben, ordnet es angestrebte Cybersicherheitsergebnisse in einer festen Gliederung und verweist für deren Umsetzung auf bestehende Kontrollkataloge~\cite[S.~i, 1]{nistcsf2}.
Diese Ergebnisse sind hierarchisch in Funktionen, Kategorien und Unterkategorien gegliedert~\cite[S.~3]{nistcsf2}.
Auf oberster Ebene unterscheidet das CSF sechs Funktionen: \emph{Govern}, \emph{Identify}, \emph{Protect}, \emph{Detect}, \emph{Respond} und \emph{Recover}~\cite[S.~3--4]{nistcsf2}.
Darunter liegen 22 Kategorien wie \emph{Risk Assessment} oder \emph{Identity Management, Authentication and Access Control}~\cite[S.~15]{nistcsf2}.

Um Risiken und Kontrollen thematisch zu ordnen, werden Kontrollen häufig zu \emph{Kontrollfamilien} zusammengefasst.
NIST nennt als Beispiel die Kontrollfamilien aus NIST SP~800~53, anhand derer sich Einträge eines Risk Registers kategorisieren lassen.
Eine einheitliche Kategorisierung erleichtert den Vergleich von Risiken~\cite[S.~12]{nistir8286}.

\section{Formalisierung}\label{sec:formal}
\begin{definition}[Risiko und Attribute]\label{def:risiko}
Ein Risiko $r$ wird durch ein Attributtupel
\[
x_r = (x_{r,1}, \ldots, x_{r,m})
\]
beschrieben.
Attribute können kategorial oder nummerisch sein oder Freitexte enthalten.
Fehlende Werte werden mit $x_{r,j} = \bot$ bezeichnet.
\end{definition}
Die Attribute entsprechen den Feldern eines Risk Registers (Abschnitt~\ref{subsec:risk_register}).
Für den Vergleich von Risiken ist zu beachten, dass insbesondere in der Entwurfsphase einzelne Werte fehlen können und Freitexte zunächst in kategoriale Attribute überführt werden müssen (Kapitel~\ref{cha:methodik}).

\begin{definition}[Kontrollfamilie]\label{def:familie}
Eine Kontrollfamilie ist eine Gruppe thematisch gleichartiger Kontrollen.
$F$ bezeichnet die Menge aller Familien.
\end{definition}
Die Familien werden in dieser Arbeit aus den Kontrollen des untersuchten Registers abgeleitet (Kapitel~\ref{cha:methodik}).

\begin{definition}[Zielgröße]\label{def:ziel}
Einem Risiko $r$ können Kontrollen zugeordnet sein. $Y_r \subseteq F$ bezeichnet die Menge der Kontrollfamilien, denen diese Kontrollen angehören.
Betrachtet werden nur Risiken mit $|Y_r| \geq 1$.
Für $n$ solcher Risiken ergibt sich die Matrix
\[
Y \in {\{0,1\}}^{n \times |F|}
\]
\end{definition}

\begin{definition}[Empfehlung]\label{def:empfehlung}
Eine Methode liefert zu einem Risiko $r$ eine nach der Einschätzung (engl.~\emph{confidence}) absteigend sortierte Liste
\[
\hat{R}_r = \big((f_1, c_1), \ldots, (f_k, c_k)\big)
\]
mit $f_i \in F$ und $c_i \in [0,1]$.
Vom Risiko $r$ selbst werden dabei nur die Attribute $x_r$ und nicht $Y_r$ verwendet.
\end{definition}
Die Einschätzung $c_i$ gibt an, wie zuverlässig die Empfehlung der Familie $f_i$ eingeschätzt wird.

\section{Ähnlichkeit gemischter Attribute}\label{sec:aehnlichkeit}
Die meisten Distanzmaße sind entweder für ausschließlich numerische oder für ausschließlich kategoriale Daten ausgelegt~\cite[S.1--2]{Velden2019}.
Die Attribute eines Risikos sind jedoch gemischt (Definition~\ref{def:risiko}).
Ein naheliegender Ausweg besteht darin, alle Attribute auf einen gemeisamen Typ zu konvertieren.
Entweder indem numerische Attribute in Klassen eingeteilt oder kategoriale Attribute numerisch kodiert werden.
Beide Wege führen jedoch zu Informationsverlust oder lassen die Dimensionen der Daten stark ansteigen~\cites[S.~82--84]{Foss2019}[S.~7]{Tarsitano2011}.
Stattdessen lassen sich Distanzmaße verwenden, die jeden Attributtypen passend behandeln und die Ergebnisse anschließend zusammenführen~\cite[S.~3]{Velden2019}.

Das bekannteste Maß dieser Art geht auf Gower zurück und gilt als Standard für gemischte Daten~\cites[s.~3]{Velden2019}[S.~89]{Foss2019}.
Für zwei Risiken $r$ und $s$ wird zunächst für jedes Attribut $j$ eine Einzeldistanz $d_j(r,s) \in [0,1]$ bestimmt, die vom Typ des Attributs abhängt.
Bei kategorialen Attirbuten is sie $0$, wenn die Werte übereinstimmen und sonst $1$.
Bei numerischen Attributen wird der Betrag der Differenz durch die Spannweite $R_j$ des Attributs geteilt:
\[
d_j(r,s) = \frac{|x_{r,j} - x_{s,j}|}{R_j}.
\]
Die Gower-Distanz ist der Mittelwert diese Einzeldistanzen,
\[
d(r,s) = \frac{1}{m} \sum_{j=1}^{m} d_j(r,s)
\]
und liegt damit zwischen $0$ und $1$~\cite[S.~4--5]{DOrazio2021}.

Ordinale Attribute, deren Werte eine Reigenfolge besitzen, werden durch ihre Rangposition ersetzt und anschließend wie numerische Attribute behandelt~\cites[S.~5]{DOrazio2021}[S.~9]{Tarsitano2011}.
Fehlt der Wert eines Attributs bei mindestens einem der beiden Risiken, wird es beim Vergleich ausgelassen.
Mit
\[
\delta_j(r,s) =
\begin{cases}
  0 & \text{falls } x_{r,j} = \bot \text{ oder } x_{s,j} = \bot,\\
  1 & \text{sonst}
\end{cases}
\]
ergibt sich die Distanz als Mittelwert über die gemeinsam vorhandenen Attribute,
\[
d(r,s) = \frac{\sum_{j=1}^{m} \delta_j(r,s)\, d_j(r,s)}{\sum_{j=1}^{m} \delta_j(r,s)}
\]
~\cites[S.~5--6]{DOrazio2021}[S.~8]{Tarsitano2011}.
Haben zwei Risiken kein gemeinsames Attribut, ist die Distanz nicht definiert~\cite[S.~6]{DOrazio2021}

Ungewichtet tragen die Attributtypen nicht gleichmäßig zur Gesamtdistanz bei.
Kategoriale Attribute erreichen bereits bei jeder Abweichung den Höchstwert $1$, numerische dagegen nur bei maximalem Abstand.
Dadurch können kategoriale Attribute die Distanz dominieren~\cites[S.~7--8]{DOrazio2021}[S.~3]{Liu2024}[S.~90]{Foss2019}.
Abhilfe schaffen Gewichte $w_j \geq 0$ je Attribut:
\[
  d(r,s) = \frac{\sum_{j=1}^{m} w_j\, \delta_j(r,s)\, d_j(r,s)}{\sum_{j=1}^{m} w_j\, \delta_j(r,s)}.
\]
Mit $w_j = 1$ für alle Attribute ergibt sich die ungewichtete Form.
Wie die Gewichte zu wählen sind, ist nicht allgemein festgelegt~\cites[S.~89]{Foss2019}[S.~3]{Velden2019}.
Sie können sich etwa an der Wichtigkeit eines Attributs oder an seinem Zusammenhang mit einer Zielgröße orientieren~\cites[S.~3--5]{Liu2024}[S.~10]{Tarsitano2011}.

\section{Attributselektion}\label{sec:selektion-grundlagen}
Nicht alle Attribute eines Risikos tragen gleich viel zur Ähnlichkeit bei, die für eine Empfehlung von Kontrollfamilien relevant ist.
Die Attributselektion bestimmt deshalb, welche Attribute wie stark berücksichtigt werden.
Ihr Ergebnis sind die Gewichte $w_j$ der Gower-Distanz (Abschnitt~\ref{sec:aehnlichkeit}).

Da einem Risiko mehrere Kontrollfamilien zugeordnet sein können, liegt ein Multi-Label-Problem vor.
Der einfachste Weg, ein solches Problem zu bearbeiten, ist \emph{Binary Relevance}.
Dabei wird die Aufgabe in unabhängige binäre Teilaufgaben zerlegt, eine je Label, und für jede Teilaufgabe wird ein eigener binärer Klassifikator trainiert~\cite[S.~191--192]{Zhang2018}.
Der Ansatz lässt sich mit jedem binären Lernverfahren kombinieren, berücksichtigt jedoch keine Zusammenhänge zwischen den Labels~\cite[S.~191--193]{Zhang2018}.

Als binärer Klassifikator dient in dieser Arbeit GRANDE (\emph{Gradient-Based Decision Tree Ensembles}).
Baumbasierte Ensembles gelten für tabellarische Daten mit gemischten Attributtypen weiterhin als besonders leistungsfähig~\cite[S.~1]{marton2024grande}.
GRANDE lernt ein Ensemble harter, achsenparalleler Entscheidungsbäume, deren Parameter nicht wie bei klassischem Boosting nacheinander, sondern gemeinsam per Gradientenabstieg optimiert werden~\cite[S.~1--3]{marton2024grande}.
Zusätzlich erhält jedes Blatt eines Baums ein eigenes Gewicht, sodass einzelne Bäume sich auf bestimmte Bereiche der Daten spezialisieren können~\cite[S.~4--5]{marton2024grande}.

Wie wichtig ein Attribut für ein trainiertes Modell ist, lässt sich mit der \emph{Permutation Importance} bestimmen.
Dazu werden die Werte des Attributs in den Validierungsdaten zufällig vertauscht und der Rückgang der Vorhersagegüte gemessen, in dieser Arbeit über die Fläche unter der ROC-Kurve~\cites[S.~1341]{Altmann2010}[S.~279]{Kaneko2022}.
Bei wenigen Datenpunkten ist das Verfahren jedoch instabil, und stark korrelierte Attribute erhalten eine zu geringe Wichtigkeit~\cite[S.~279]{Kaneko2022}.
Entscheidend ist, dass die Selektion ausschließlich auf den Trainingsdaten erfolgt.
Fließen die Testdaten in die Auswahl ein, fällt die anschließende Bewertung zu optimistisch aus~\cites[S.~6562]{Ambroise2002}[S.~1--2]{Varma2006}.
\section{Nachbarschafts- und clusterbasierte Empfehlung}\label{sec:empfehlung-grundlagen}
\todo{kNN und ML-kNN~\cite{zhang2007mlknn}, Clustering kategorialer Daten~\cite{huang1998kmodes}, Bezug zu CBR~\cite{aamodt1994cbr}.}


\section{Bewertung von Empfehlungen}\label{sec:metriken}
Für ein Test-Risiko $r$ mit den Top-$k$-Familien $T_r$ gilt
\begin{align}
  \mathrm{P@}k(r) &= \frac{|T_r \cap Y_r|}{k}, &
  \mathrm{R@}k(r) &= \frac{|T_r \cap Y_r|}{|Y_r|}.
\end{align}
Die Kalibrierung der Confidence wird über den Expected Calibration Error mit
$B$ gleich breiten Bins gemessen:
\begin{equation}
  \mathrm{ECE} = \sum_{b=1}^{B} \frac{|B_b|}{N}\,\Big|\overline{c}(B_b) - \overline{h}(B_b)\Big|
\end{equation}
\todo{Metriken erläutern, $k=3$ begründen, Literatur zur Kalibrierung ergänzen.}
